% ============================================================
%  SEZIONE 3 — Funzionalità
% ============================================================
\section{Funzionalità}

\subsection{Descrizione della funzionalità selezionata}

% Scegliete UNA funzionalità significativa da descrivere nel dettaglio.
[Descrivere qui la funzionalità scelta, il suo scopo e il contesto d'uso all'interno dell'applicazione.]

\subsection{Interazione utente}

Dal punto di vista dell'utente, la funzionalità si articola nei seguenti passi:

\begin{enumerate}
  \item Passo 1: descrivere l'azione dell'utente.
  \item Passo 2: descrivere la risposta del sistema.
  \item Passo 3: eventuale feedback o esito finale.
\end{enumerate}

\subsection{Flusso interno al sistema}

La Figura~\ref{fig:sequenza} mostra il diagramma di sequenza relativo all'esecuzione della funzionalità, evidenziando i componenti coinvolti e i messaggi scambiati.

\begin{figure}[H]
  \centering
  % --- Diagramma di sequenza con TikZ ---
  \begin{tikzpicture}[
      actor/.style  = {draw, thick, fill=primary!10, minimum width=2.2cm,
                       minimum height=0.7cm, rounded corners=3pt,
                       font=\small\bfseries, align=center},
      lifeline/.style = {draw, dashed, thick, color=secondary!60},
      msg/.style    = {-Stealth, thick, color=primary!80},
      retmsg/.style = {-Stealth, dashed, thick, color=secondary},
      lbl/.style    = {font=\scriptsize, fill=white, inner sep=1pt},
    ]

    % --- Attori (asse X) ---
    \def\xBrowser{0}
    \def\xCtrl{3.5}
    \def\xModel{7}
    \def\xDB{10.5}

    \node[actor] (A) at (\xBrowser, 0) {Browser};
    \node[actor] (B) at (\xCtrl,    0) {Controller};
    \node[actor] (C) at (\xModel,   0) {Model};
    \node[actor] (D) at (\xDB,      0) {Database};

    % --- Lifelines ---
    \foreach \x in {\xBrowser, \xCtrl, \xModel, \xDB}{
      \draw[lifeline] (\x, -0.35) -- (\x, -5.8);
    }

    % --- Messaggi (asse Y scende) ---
    % 1. Request
    \draw[msg] (\xBrowser,-0.8) -- node[lbl,above]{1. HTTP Request} (\xCtrl,-0.8);
    % 2. validate & fetch
    \draw[msg] (\xCtrl,-1.4) -- node[lbl,above]{2. [validate] fetch()} (\xModel,-1.4);
    % 3. query
    \draw[msg] (\xModel,-2.0) -- node[lbl,above]{3. SQL query} (\xDB,-2.0);
    % 4. result set
    \draw[retmsg] (\xDB,-2.6) -- node[lbl,above]{4. result set} (\xModel,-2.6);
    % 5. domain object
    \draw[retmsg] (\xModel,-3.2) -- node[lbl,above]{5. domain object} (\xCtrl,-3.2);
    % 6. render view
    \draw[msg] (\xCtrl,-3.8) -- node[lbl,above]{6. render(view, data)} (\xBrowser+0.3,-3.8);
    % piccola freccia interna al controller per indicare passaggio alla view
    \draw[msg] (\xCtrl,-4.4) -- node[lbl,above]{7. View populated} (\xCtrl+0.01,-4.4);
    % 8. HTTP Response
    \draw[retmsg] (\xCtrl,-5.2) -- node[lbl,above]{8. HTTP Response (HTML)} (\xBrowser,-5.2);

  \end{tikzpicture}
  \caption{Diagramma di sequenza della funzionalità selezionata.}
  \label{fig:sequenza}
\end{figure}

\subsection{Funzionalità non implementate}

% Descrivere eventuali funzionalità previste ma non realizzate, con motivazione.
[Elencare e motivare le funzionalità che non è stato possibile implementare, specificando se si tratta di vincoli temporali, tecnici o di altra natura.]
